\begin{prob}
    Determine whether each piece of code is correct or
    incorrect. If it is correct, say so (you do not need to show work). But if it
    is incorrect, do the following:
    
    \begin{enumerate}
        \item Say what will go wrong (e.g., infinite recursion, incorrect
            answer) and give an example of an input that will demonstrate the
            error.
        \item Fix the code so that the function works properly by
            \textbf{either} 1) adding one or more lines of code; 2) changing
            one (and only one) line in the code given.  (You won't need to add
            code \emph{and} change an existing line. You're not allowed to
            remove a line of code.)
    \end{enumerate}

    Hint: use the ideas behind inductive proofs to try and argue to yourself that a piece
    of code is correct.

    \begin{subprobset}
        \begin{subprob}
            \inputminted{python}{\thisdir/include/part_a.py}
            Recall that in Python if \python{s} and \python{t} are strings, then \python{s + t}
            is the concatenation of \python{s} and \python{t}.
            \begin{soln}
                This is missing a base case for when the string is empty. If the string is
                empty, we will get an index error when we try to access \python{s[0]}.

                The fix is to add a base case for when the input array is of size 0.
                \inputminted{python}{\thisdir/include-solution/part_a.py}
            \end{soln}
        \end{subprob}

        \begin{subprob}
            In this part, adopt the convention that the product of the numbers in an empty
            list is one 
            (why? since $b^0 = 1$, whatever $b$ is, by convention).
            \inputminted{python}{\thisdir/include/part_b.py}
            \begin{soln}
                This is correct.
            \end{soln}
        \end{subprob}

        \begin{subprob}
            You may assume that \python{start} and \python{stop} are valid indices.
            \inputminted{python}{\thisdir/include/part_c.py}
            \begin{soln}
                The arguments to the recursive call are wrong: they should sort the middle part
                of the list (excluding the first and last element), but they include the first
                element. This causes incorrect output. For instance, \python{inplace_reverse([4,5,2,3], 0, 4)} results in
                \python{[3,5,2,4]}.

                The fix is to change the argument:
                \inputminted{python}{\thisdir/include-solution/part_c.py}
            \end{soln}
        \end{subprob}

        \begin{subprob}
            \inputminted{python}{\thisdir/include/part_d.py}
            \begin{soln}
                This works correctly.
            \end{soln}
        \end{subprob}
    \end{subprobset}
\end{prob}
